% ############################################################
% VERDICT AFTER HOSTILE REVIEW: DO NOT ADD THIS SECTION.
% The paper already states its thesis, as an axiom, better.
%
% mechviews_v19.tex:386-390, coherence condition (ii):
%   "Evidence tracks identity: measurements should not distinguish
%    objects that the identity criterion declares equivalent. (Evidence
%    that discriminates within an identity class is tracking a
%    finer-grained ontology than the one declared.)"
% That IS over-resolution relative to a declared identity criterion --
% the three-place claim this section announces as new. VERIFIED by
% reading the lines.
% The free-parameter half is at :1062-1071 (analyst-choice family,
% "an artifact of a choice made before any measurement was taken"), and
% :1299 already runs the sharpens-vs-over-resolves discriminator on the
% SAE-width case. The Stratified view at :849-853 already makes identity
% resolution-parameterized and names under-resolution as a failure mode.
%
% ALSO FALSE, and I told Elliot the opposite earlier -- CORRECTED:
% Cao & Yamins' L5PC 7-layer MLP is NOT an over-resolution case.
% cao_yamins_part1.txt:1243-1246 -- it fails 3M++ "because it does not
% satisfy a model-to-mechanism mapping ... the components of the MLP do
% not map to real-world components." That is NON-CORRESPONDENCE, a
% different failure. There is still no worked over-resolution case
% anywhere in the sources.
%
% GEIGER OWNS THE RELATIVIZATION. Remark 32 (:1022): "An alignment
% induces a unique partial function omega^pi that maps from low-level
% hard interventions to high-level hard interventions" -- intervention
% distinguishability indexed to the declared alignment. Definition 38
% (value merge) IS a declared identity criterion. Sec 3.6 catalogs PCA,
% SAEs, probes and DAS as four coexisting featurizations of one vector,
% i.e. several criteria at once. DAS SEARCHES alignments, so "Geiger
% fixes an alignment" is false about the cited paper.
%
% FANG & ZHANG (Erkenntnis) ALREADY RELATIVIZES, AND HAS THEOREMS:
% :274-277 the cause-variable is a partition induced by an equivalence
% relation, relative to Y; footnote 12 (:656-659) relative to a chosen
% possibility space; footnote 15 (:868-869) base level "may involve some
% convention or decision". Theorem 1 (:498-499) "determinate causes sink
% down" is the over-resolution case with a proof; non-preservation under
% coarse-graining (:731-735) is under-resolution.
% AND IT HAS THE COUNTER-THESIS: :432-434, "there is a UNIQUE
% proportional cause-variable ... for a given effect-variable."
%
% WHAT TO DO INSTEAD:
%  - Move the matched/under/over vocabulary into Sec 4 as three
%    sentences, cited to Yablo and Cao & Yamins. It is useful shared
%    wording. It cannot carry a subsection.
%  - Fold the sharpens-vs-over-resolves criteria into Sec 5 next to
%    :1299, where the SAE case already runs them.
%  - The ONE live claim: a single interpretability explanandum supports
%    four inequivalent identity criteria (component set, role, subspace,
%    gauge orbit). Fang & Zhang's uniqueness theorem is the named
%    opponent. This must be DEMONSTRATED on a case where the four
%    criteria disagree about one target. Asserted, it is worthless.
%  - The path-patching example below is WRONG: mechviews_v19.tex:1106
%    classifies path patching as Object/component-overlap/directed-graph,
%    and :1474 lists it as legitimate Object evidence. Edges are objects
%    of the declared formalism, so the routes are not "surplus".
% ############################################################

% Resolution matching — the section that must exist before any positioning.
% Candidate home: §4 (Atlas), as a subsection preceding the nine view entries,
% so every view entry can use the vocabulary.
%
% WRITTEN AGAINST THE HOSTILE REVIEW. The earlier positioning draft claimed
% over-resolution as new. It is not: Cao & Yamins have it with a worked case
% (cao_yamins_part1.txt:1524-1530), and Geiger et al. have both directions
% (alignment induces an intervention map; pi undefined on unrealized values).
% This section concedes both and locates the contribution elsewhere — in the
% choice of identity criterion, which is the free parameter neither fixes.
%
% QUOTE PROVENANCE, all verified in local PDFs:
%   cao_yamins_part1.txt:1524-1530  "the complexity of the description blows up"
%   cao_yamins_part1.txt:742-744    runnability forces detail / "armchair"
%   geiger_2023_causal_abstraction.txt:1-10  JMLR 26 (2025), arXiv v4
% BIBITEMS REQUIRED: caoyamins2024a, geiger2025, yablo1992
% NOTE: cite Geiger as 2025 (JMLR 26), NOT 2023. Check collision with the
% existing geiger2024 bibitem used at v19:693.

\subsection{Matching an Intervention to a Mechanism}
\label{sec:resolution}

A claim about a mechanism rests on a measurement, and a measurement can be too
blunt or too sharp for the thing it measures.  Call the finest distinction an
intervention can draw its \emph{resolving power}, and call the distinctions a
mechanism instantiates its \emph{grain}.  Three states follow.  A report is
\emph{matched} when resolving power meets grain.  It is \emph{under-resolved}
when the intervention is blunter than the mechanism, so real structure goes
unreported.  It is \emph{over-resolved} when the intervention draws distinctions
the mechanism does not carry, so the surplus structure in the report is an
artifact of the probe.

Both failures are known.  \citet{caoyamins2024a} match the grain of a model to
the grain of its explanandum and supply a worked case of excess resolution: a
seven-layer network fitted to a single cortical cell, where ``the complexity of
the description blows up as its fidelity to the physical mechanism falls
apart.''  They also supply the standing caution against assuming detail is
discardable, since runnability ``can end up forcing the preservation of many
internal details'' beyond what one would guess ``from the armchair.''
\citet{geiger2025} carry the same structure into interpretability.  An alignment
between a high-level model and a network induces a map from low-level
interventions to high-level ones, so coarsening the high-level variables
coarsens which interventions remain distinguishable, and interchange-intervention
accuracy grades how well the pairing holds.  The vocabulary below renames
neither result.

What both proposals hold fixed is the interesting part.  Cao and Yamins fix the
explanandum, then ask which model grain suits it.  Geiger et al.\ fix an
alignment, then ask whether the network realizes it.  Once either is fixed the
grain question is well posed, and each answers it.  Choosing what to fix is the
prior question, and it is a choice among identity criteria: an alignment
declares which states count as the same mechanism before any intervention runs.
A network admits several such declarations at once.  The same head participates
in a component set, a functional role, a subspace, and a gauge orbit, and each
carries its own answer to what would count as the same mechanism appearing
twice.

Resolution matching is therefore a three-place relation that the literature
treats as two-place.  An intervention resolves finely or coarsely
\emph{relative to a declared identity criterion}, because the criterion fixes
which distinctions count as distinctions at all.  Path patching separates the
routes by which a component's effect travels.  Under object identity those
routes are one mechanism, and the separation is surplus.  Under role identity
they may be two, and the same measurement is exactly matched.  Neither reading
is a better measurement.  They differ in what was declared before the
measurement was taken, and a report that omits the declaration leaves its own
resolution unassessable.

This yields four outcomes when two interventions disagree about one claim.  They
\emph{agree} when their reports are consistent at some shared level.  They
\emph{disagree} when the reports are jointly inconsistent at every level.  One
\emph{sharpens} the other when the finer intervention reports additional
structure that survives verification.  One \emph{over-resolves} when the
additional structure fails it.  Separating the last two takes evidence this
framework does not supply: reproducibility of the additional structure across
seeds, its predictive purchase on held-out behavior, and its coherence with
reports at other levels of description.  Those are audit questions, and
\citet{tower2026mechvalidity} supplies the audits.  The contribution here is the
statement of what the audit is being run on, which requires a declared
criterion, and which is why the atlas comes first.
